% Generated draft: requires mathematical and source review.
\documentclass[11pt]{article}
\usepackage[margin=25mm]{geometry}
\usepackage{fontspec}
\setmainfont{Latin Modern Roman}
\usepackage{amsmath,amssymb,mathtools,unicode-math}
\setmathfont{Latin Modern Math}
\usepackage{enumitem,needspace,xurl,hyperref}
\setlength{\parindent}{0pt}
\setlength{\parskip}{5pt}
\setlength{\emergencystretch}{3em}
\newcommand{\N}{\mathbb{N}}
\newcommand{\Z}{\mathbb{Z}}
\newcommand{\Q}{\mathbb{Q}}
\newcommand{\R}{\mathbb{R}}
\newcommand{\C}{\mathbb{C}}
\newcommand{\F}{\mathbb{F}}
\newcommand{\A}{\mathbb{A}}
\newcommand{\PP}{\mathbb{P}}
\newcommand{\E}{\mathbb{E}}
\newcommand{\eps}{\varepsilon}
\newcommand{\dd}{\,\mathrm{d}}
\begin{document}
\section*{UnsolvedMath ID: 10900005}
\textbf{Source problem code:} catalogue entry
\subsection{Definitions and mathematical statement}
Let $N$ be a closed $3$-manifold whose JSJ decomposition consists entirely of hyperbolic pieces. In the standard setting, the decomposition cuts along JSJ tori, and each resulting piece has interior admitting a complete finite-volume metric of constant sectional curvature $-1$.

Danciger's question, item 1.5, is whether every such $N$ admits a convex real projective structure. In the customary properly convex interpretation, this asks for a properly convex open domain $\Omega\subset\mathbb{RP}^3$ and a discrete group $\Gamma\subset\operatorname{PGL}_4(\R)$ preserving $\Omega$, acting freely and properly discontinuously with compact quotient, such that
\[
 N\cong\Omega/\Gamma.
\]
Proper convexity means that $\Omega$ is convex and bounded in some affine chart. This must not be strengthened to strict convexity: the source asks only for convexity. The excerpt does not specify whether proper convexity is part of its convention, nor explicitly state the usual connectedness, orientability and irreducibility hypotheses under which its JSJ terminology is being used. These conventions must be checked rather than extending the assertion to unrelated reducible cases.

The source identifies this as the more general version of item 1.4, whose statement is absent here. The registry reports gluing constructions in particular cases and says the general all-hyperbolic-JSJ statement remains unproved; this is supplied progress information.
\subsection{Short English statement}
Must a closed $3$-manifold built entirely from hyperbolic JSJ pieces carry a convex real projective structure? The exact convexity and JSJ conventions need source confirmation.
\subsection{Sources}
\begin{itemize}
\item Local registry: \url{data/unsolvedmath/problems.json}, permanent UnsolvedMath ID $10900005$, exact source code \texttt{catalogue entry}, attributed to Danciger.
\item Delp, Hoffoss and Manning, \emph{Problems In Groups, Geometry, and Three-Manifolds} (2015), subsection 1.5, PDF page 1: \url{https://arxiv.org/pdf/1512.04620#page=1}; abstract record: \url{https://arxiv.org/abs/1512.04620}. Supplied access date: 2026-07-29; source-tex extraction checked against rendered PDF text by the source record.
\item The registry mentions gluing technology, doubles, and some graph manifolds, without precise additional references. Graph-manifold examples are related context and are not thereby examples satisfying the all-hyperbolic-JSJ hypothesis.
\end{itemize}
\end{document}
